\scanpage{70}
\mspage{\handwritten{- 5 -}}

p.11, li.7 du bas : $K^{\bullet}(T)$ et non $K_{\bullet}(T)$.

p.11, li.18 : il devrait y avoir une référence dans Exp.V pour cette assertion.

p.11, li.11 du bas : référer à (3.1) pour la structure de module envisagée.

p.11, li.9 du bas : ``ci-dessus'' et non ``ci-dessous''. Apparemment les définitions garderaient un sens en supposant seulement $C$ connexe, ou est-ce que je me canule ?

p.12,li.6 : ``on se ramène immédiatement ...'' : bref, on utilise la compatibilité de la spécialisation au changement de trait\handwritten{(7\unsure{.|-}15)}.

p.12, li.6 du bas : ``donc il existe (cf. EGA IV 2.8) ...''

p.12, li.3 du bas : ``on peut choisir pour $C$ un schéma régulier'', mais plus nécessairement plat sur $S$.

p.13, li.2 : ``il est bien connu (SGA ou EGA ? ) ...''

\handwritten{p.13, li.10 du bas : ($\overline{\underline{\mathrm{IV}}}$ 2.12) au lieu de ($\overline{\underline{\mathrm{III}}}$ 2.12).}

p.13, li.3 du bas : il semble qu'on utilise le fait déjà cité(sans référence) que par $x_{os}$ et $x_{\unsure{l|1}s}$ passe une courbe connexe.

p.14, li.10 : ``propre et lisse'' ne suffit-il pas ?

p.15, li.10 : ``donne naissance à '' plutôt que ``provient de''.

p.16, li.14 : au lieu de ``par Krull-Akizuki'' dire plutôt ``comme $S$ est hensélien''.

p.16, li.16 : $X$ n'étant pas supposé plat, le fait que les fibres soient géométriquement régulières n'implique pas, il me semble, que $X$ soit lisse.

\handwritten{p.16, li.8 du bas : ``on trouve''\unsure{...|.} ajouter ``(cf. 7.16)''.}

p.17, 7.13.6 : les points importants, à savoir :
\begingroup
\makeatletter\tagsleft@true\makeatother
\[
\handwritten{\psi}_i(c^i(x)) = (-1)^{i-1}(i-1)!\;c^i_{\handwritten{\chi}}(x)\quad,\tag*{(i)}
\]
\[
\handwritten{\psi}_{i+j}(c^i(x)c^j(y)) = (-1)^{i+j-1}(i+j-1)!c^i_{\handwritten{\chi}}(x)c^j_{\handwritten{\chi}}(y)\quad,\tag*{(ii)}
\]
\endgroup
pour lesquels on peut référer à (V \handwritten{6}.\handwritten{5}.\handwritten{3} et \handwritten{6}.\handwritten{6}.\handwritten{4}), ne sont pas
% archive sheet number: 69
